%% ============================================================
%%  Section: Identification Strategy
%%  Depositor Discipline in Modern Banking (2026-04-04 draft)
%%
%%  Cross-references:
%%    eq:dayone, eq:treatment, eq:binary  (sec:background)
%%    tab:descriptive_sumstats_cecl, tab:cecl_predictors_cross_section
%%    fig:cecl_equity_distribution
%%    sec:results, sec:results:heterogeneity
%%
%%  New citations suggested (see end of file):
%%    callaway2021difference, sun2021estimating, roth2023whats
%% ============================================================

\section{Identification Strategy}
[@@@ overall this section is too long and don't pitch the natural experiment as a silver bullet. be less confident, but communicate that it is a good/resonable natural experiment that gets us a causal interpretation. Reduce repeating stuff from previous sections, particuraly from institutional background]
\label{sec:identification}
[@@@ less --- emdashes]

Our goal is to identify the causal effect of a credit-risk disclosure on
depositor behavior.  The empirical obstacle is familiar: riskier banks carry
fewer uninsured deposits in equilibrium, so cross-sectional comparisons of
deposit structure on bank risk conflate depositors' information responses with
pre-existing sorting and equilibrium co-determination
\citep{covitz2004market, chen2022bank}.  We resolve this problem by exploiting
the CECL Day-One adjustment as an instrument[@@@ natural experiment to be consistent with previous sections] that delivers a precisely
identified information shock[@@@ too strong claim] --- predetermined in both timing and magnitude,
separable from concurrent changes in bank cash flows, and concentrated enough
in calendar time to be isolated from confounding economic trends.
This section formalizes the identification argument, develops the estimating
equations, and documents the information content of the CECL disclosure using
equity market and CDS evidence from the large-bank adoption cohort. [@@@ this paragaph is too strong, overselling]

\subsection{Empirical Design}
\label{sec:identification:did}
[@@@ remove paragraph headers like in the previous sections]
\paragraph{Static difference-in-differences.}
Our baseline specification is a difference-in-differences estimator that
exploits cross-sectional variation in CECL exposure and the common
regulatory adoption date:
\begin{equation}
  y_{it}
  = \alpha_i + \gamma_t
  + \beta \cdot \mathrm{Post}_{it} \times \mathrm{CECL}_i
  + \mathbf{X}_{it-1}'\boldsymbol{\delta}
  + \varepsilon_{it},
  \label{eq:did}
\end{equation}
where $y_{it}$ is the outcome for bank $i$ in quarter $t$; $\alpha_i$ and
$\gamma_t$ are bank and calendar-quarter fixed effects; $\mathrm{Post}_{it}$
equals one for all quarters at or after bank $i$'s CECL adoption date;
$\mathrm{CECL}_i$ is either the continuous \texttt{CECL Adj} measure or the
binary \texttt{High CECL} indicator, both defined in
Section~\ref{sec:background}; and $\mathbf{X}_{it-1}$ is a vector of lagged
controls (log assets, equity-to-assets ratio, interest expense-to-assets).
Bank fixed effects absorb all time-invariant differences across institutions
--- balance-sheet composition, business model, geographic market, charter type
--- so that $\hat{\beta}$ is identified from within-bank changes in outcomes
over time.  Quarter fixed effects absorb any aggregate macroeconomic force
common to all banks in a given quarter, including the Federal Reserve's
tightening cycle, aggregate deposit normalization, and any nationwide
reaction to the 2023 banking stress.

\paragraph{Event-study specification.}
To evaluate the parallel-trends assumption and trace the dynamics of depositor
responses, we estimate:
\begin{equation}
  y_{it}
  = \alpha_i + \gamma_t
  + \sum_{\substack{k=-10 \\ k \neq -1}}^{10}
    \beta_k \cdot \mathbf{1}\!\left\{\tau_{it} = k\right\}
    \times \mathrm{CECL}_i
  + \mathbf{X}_{it-1}'\boldsymbol{\delta}
  + \varepsilon_{it},
  \label{eq:eventstudy}
\end{equation}
where $\tau_{it} = t - t_i^{*}$ is event time in quarters relative to bank
$i$'s CECL adoption date, and $k = -1$ is the omitted reference period.
Coefficients $\{\hat{\beta}_k\}_{k < 0}$ test for pre-adoption differential
trends between high- and low-CECL banks; coefficients $\{\hat{\beta}_k\}_{k
\geq 0}$ trace the cumulative post-adoption effect.  The identifying assumption
requires that the pre-adoption coefficients be jointly zero.  A pattern of flat pre-trends followed by a post-adoption break
supports the interpretation that observed differences are caused by the CECL
disclosure rather than pre-existing divergence.

\paragraph{Staggered adoption.}
Ninety-two percent of sample banks adopted CECL in 2023Q1, with the remainder
spread across the five adjacent quarters in the 2022Q3--2023Q4 window.  This
near-uniform adoption calendar implies that the biases identified by
\citet{callaway2021difference}, \citet{sun2021estimating}, and
\citet{dechaisemartin2020two} for staggered difference-in-differences designs
--- which arise when later-treated units serve as implicit controls for
earlier-treated ones --- are quantitatively minor in our setting.[@@@ these are commonly known papers, dont cite. Don't emphasize too much the uniform adoption calendar]  We
nonetheless confirm this by reporting stacked-cohort estimates that
restrict control banks to those adopting in the same quarter, and by
implementing the \citet{callaway2021difference} group-time average treatment
effect estimator.  Both robustness checks yield results nearly identical to
the pooled specification in \eqref{eq:did} [@@@ we don't do this in the code. (last two sentences.). Don't add things not in the code].

\subsection{Treatment Exogeneity}
\label{sec:identification:predetermination}

For $\hat{\beta}$ in \eqref{eq:did} to have a causal interpretation,
the CECL treatment must be uncorrelated with potential outcomes absent the
disclosure.  This requires that neither the timing of the shock nor its
magnitude was determined by --- or correlated with --- the deposit dynamics we
study.  Both conditions hold by institutional design.

\paragraph{Timing is regulatory, not managerial.}
The 2023Q1 adoption date was set by FASB's ASU 2016-13 and enforced by the
federal banking agencies through their regulatory capital rules
\citep{FASB2016}.  Institutions subject to the community-bank rollout had
no ability to accelerate adoption to preempt scrutiny, delay it to conceal an
unfavorable charge, or time it to coincide with favorable deposit market
conditions.  The adoption date is therefore orthogonal to any bank-specific
information about contemporaneous deposit flows.

\paragraph{Magnitude reflects historical lending, not current conditions.}
The size of $\Delta_i$ in equation~\eqref{eq:dayone} is determined by the
cumulative composition of the loan portfolio as it stood on the adoption date,
itself the product of lending decisions made years or decades prior.  The
CECL charge is large when a bank's portfolio contains loans with long remaining
maturities, weak borrower credit quality, or thin collateral coverage ---
characteristics that were fixed by prior underwriting decisions, not by
2023 deposit market conditions.  Some banks may have gradually tightened
underwriting after 2016 in anticipation of CECL, but this response
\textit{compresses} rather than eliminates cross-sectional variation in
$\Delta_i$, and the residual variation attributable to pre-2016 portfolio
composition is precisely what identifies our estimates.  The
cross-sectional regression in Table~\ref{tab:cecl_predictors_cross_section}
confirms this: the CECL adjustment is predicted by NPL ratios and historical
allowance levels --- credit-risk fundamentals --- with a maximum adjusted
$R^2$ of 1.15 percent across all specifications.  Its near-zero correlation
with deposit structure (Table~\ref{tab:descriptive_sumstats_cecl}) confirms
that treatment intensity is orthogonal to the pre-adoption level of the
outcomes we study. [@@@ too long. we cannot confidently call these orthogonal. be less confident. a referee can always comeup with pushback]

\paragraph{Phase-in election as a mechanism test.}
Banking regulators permitted community banks to elect a three-year phase-in
of the Day-One charge against regulatory capital ratios.  Banks electing the
phase-in (\texttt{phase\_in}$_i = 1$) faced an identical informational
disclosure to non-electing banks --- the full CECL reserve and retained-earnings
charge appear in the Call Report regardless of phase-in status --- but
experienced a more gradual mechanical impact on their regulatory capital
ratios over 2023--2025.  This separation between information and capital
provides a direct test of the mechanism.  If depositor responses are driven
by the informational content of the disclosure, they should be present
regardless of phase-in election.  If they are driven by the capital reduction
per se, they should be attenuated or absent for phase-in banks whose regulatory
ratios are insulated.  Section~\ref{sec:results:heterogeneity} documents that
both high- and low-phase-in banks respond, with the triple-interaction
coefficient statistically indistinguishable between the two groups, confirming
that information rather than capital mechanics drives the effect. [@@@ not sure if this information is needed for this paper. take another look.]

\subsection{Information Content, Not Cash Flows}
\label{sec:identification:mechanism}

A central advantage of our design relative to the depositor-discipline
literature is the separation of information from concurrent changes in bank
financial condition.  Most prior tests exploit realized distress events ---
rating downgrades, dividend cuts, NPL spikes, near-failure episodes --- in which
new information arrives simultaneously with deterioration in the bank's actual
cash flows, collateral values, or liquidity position
\citep{flannery1996evidence, martinez2001depositors, iyer2016tale}.  In such
settings, the depositor response to the information \textit{per se} is
inseparable from the response to the contemporaneous financial deterioration.
The CECL Day-One adjustment resolves this confound directly.  As emphasized in
Section~\ref{sec:background}, the adjustment records a pure accounting charge:
it does not trigger a cash transfer, alter borrower repayment terms, reduce
bank liquidity, or change the risk of the underlying loans.  A depositor who
withdraws funds or demands a higher rate after observing a large CECL
adjustment is responding to newly disclosed information about the credit-risk
composition of the bank's loan book --- not to a concurrent worsening of the
bank's actual financial condition.

This design also dominates tests based on regulatory actions and supervisory
ratings changes, where the information event is confounded by the direct
behavioral constraints imposed by regulators.  And it dominates tests based
on voluntary disclosure events, where the choice to disclose is itself
informative and endogenous to management strategy.  The CECL disclosure is
mandatory, quantified, audited, and timed by regulation --- properties that
make it an unusually clean instrument for studying how depositors process
financial information.

\subsection{Threats to Identification} [@@@ this should be much shorter. one paragraph with few sentences]
\label{sec:identification:threats}

\paragraph{Anticipation.}
If sophisticated depositors had already priced in the CECL adjustment before
the official adoption date, the post-adoption treatment effect would be
attenuated toward zero.  We address this concern with two independent pieces
of evidence.  First, the pre-trend coefficients $\{\hat{\beta}_k\}_{k < 0}$
in our event-study specifications are uniformly flat and statistically
indistinguishable from zero across the full ten-quarter pre-adoption window
for all outcome variables.  Second, and independently, equity markets and
CDS spreads at large SEC-filing banks --- which have stronger incentives and
greater capacity to anticipate accounting disclosures than community bank
depositors --- also show flat pre-adoption responses across the full nine-month
pre-event window, with sharp breaks only at the adoption month.  This evidence,
detailed in Section~\ref{sec:identification:signal}, establishes that even
the most sophisticated financial market participants did not price in the CECL
adjustment before it was disclosed.  The low $R^2$ in the cross-sectional
CECL predictor regressions (Table~\ref{tab:cecl_predictors_cross_section})
provides a further explanation: the CECL adjustment is not sufficiently
predictable from public information for depositors to have fronted the signal. [@@@ don't need this paragraph]

\paragraph{The 2023 banking stress episode.}
A direct threat to identification is that 2023Q1 --- the primary CECL adoption
quarter --- coincides with the failures of Silicon Valley Bank and Signature Bank
in March 2023, which triggered widespread attention to uninsured deposit risk
across the banking system.  If depositors at high-CECL banks responded to
the aggregate stress episode rather than to CECL-specific information, the
treatment effect would be confounded.  Three features of our design rule this
out.  First, the 2023 stress was driven by unrealized securities losses and
duration mismatches at interest-rate-sensitive institutions, not by credit
risk.  Table~\ref{tab:cecl_predictors_cross_section} shows that unrealized
securities losses have essentially zero correlation with the CECL treatment
variable (coefficient 0.0008, $t < 1$), implying that the banks most exposed
to the 2023 stress channel are not disproportionately represented in our
High-CECL group.  Second, calendar-quarter fixed effects absorb the aggregate
shock to uninsured deposit flows common to all community banks in 2023Q1,
including any systemic response to the SVB episode; the treatment effect is
identified from the \textit{differential} response at high-CECL relative to
low-CECL banks within that quarter, which is uncorrelated with the common
stress trigger.  Third, the community banks in our sample are substantially
insulated from the 2023 stress: all have assets below \$10~billion, carry
primarily traditional loan-funded balance sheets, and lack the concentrated
technology-sector uninsured deposit bases that made SVB uniquely fragile
\citep{jiang2024monetary}.

\paragraph{Monetary tightening and deposit repricing.}
The Federal Reserve raised the federal funds rate by 525 basis points between
March 2022 and July 2023, causing significant deposit repricing across the
banking system.  If high-CECL banks were differentially exposed to the rate
cycle --- for example, through shorter-duration liabilities or greater reliance
on rate-sensitive depositors --- their deposit dynamics would reflect rate
pass-through rather than the CECL information shock.  Two features of our
specification address this concern.  The lagged interest expense-to-assets
control absorbs pre-existing differences in deposit pricing strategy across
institutions.  Calendar-quarter fixed effects absorb aggregate time-series
variation in rates and deposit flows.  Identification comes from the
cross-sectional interaction between rate-cycle timing and CECL exposure;
because the CECL treatment is orthogonal to pre-adoption deposit structure
(Table~\ref{tab:descriptive_sumstats_cecl}), differential rate sensitivity
is not a channel through which CECL exposure could spuriously correlate
with post-adoption deposit flows.

\paragraph{General equilibrium and spillovers.}
If uninsured deposits that leave high-CECL banks flow to low-CECL banks, the
control group's deposit outcomes would be positively contaminated, downward-biasing
the estimated treatment effect.  This concern is most relevant if the community
bank deposit market is localized, so that high- and low-CECL banks in the same
geographic market directly compete for the same depositor pools.  While we
cannot rule out this attenuation entirely, its direction is conservative: any
spillovers would bias our estimates toward zero, making it harder rather than
easier to find evidence of depositor discipline. [@@@ don't need this]

\subsection{Signal Validation: Market Responses at Large-Bank Adoption} [@@@ again too long]
\label{sec:identification:signal}

A necessary condition for our causal interpretation is that the CECL adjustment
contains information about bank credit quality that was not already reflected
in prices.  If the adjustment merely relabeled already-known risk --- or if
depositors believed it did --- no rational monitor should revise her behavior
in response to it.  We test this condition using the cohort of large, publicly
traded banks that adopted CECL in 2020Q1, exploiting the availability of
high-frequency market data unavailable for community banks.

\paragraph{Why the 2020 SEC-filer cohort.}
Large SEC-filing banks were required to adopt CECL beginning with fiscal years
starting after December~15, 2019.  For calendar-year filers, this means
2020Q1.  Because equity prices and CDS spreads are observable around this
adoption date at monthly frequency, we can test directly whether sophisticated
financial market participants --- whose pricing reflects private research,
analyst forecasts, and regulatory filings --- revised their assessments at the
moment of CECL disclosure.  The 2020 cohort cannot be used for our main
depositor-discipline tests because the adoption date coincides exactly with
the onset of the COVID-19 pandemic, which overwhelmed any deposit-market
signal with aggregate liquidity demand and emergency policy interventions.
We use this cohort solely to establish information content, carrying that
result forward as evidence that community bank uninsured depositors in 2023
were responding to genuine new information rather than a mechanical accounting
change.

\paragraph{Event-study estimates: equity markets.}
Figure~\ref{fig:es_large_banks_market_cds} Panel~A reports event-study
estimates of the differential change in normalized market capitalization for
SEC-filer banks around their CECL adoption date, using the continuous
\texttt{cecl\_equity} ratio as treatment intensity, with bank and
calendar-month fixed effects.  The pre-adoption coefficients across all nine
lead months are flat and statistically indistinguishable from zero, ruling out
anticipation by equity investors --- who had every incentive and ability to have
already priced bank credit risk from prior disclosures.  At the adoption month
(month~0), the coefficient drops sharply to $-0.947$ ($p = 0.006$), indicating
that a one-percentage-point higher CECL-to-equity ratio is associated with a
0.95-unit decline in normalized market capitalization in the adoption month
itself.  The effect deepens over the following two months, reaching a trough of
$-2.771$ ($p < 0.001$) at month~2 --- implying that a ten-percentage-point
higher CECL adjustment is associated with roughly 28~percent lower market
capitalization relative to the adoption month, after removing common monthly
shocks and bank fixed effects.  The effect partially reverts over months~3
through~9 but remains negative and economically substantial through month~6.
This pattern is inconsistent with a purely cosmetic accounting reclassification.

\paragraph{Event-study estimates: credit default swaps.}
Figure~\ref{fig:es_large_banks_market_cds} Panel~B reports the analogous
event-study using monthly CDS spreads for the 103 SEC-filing banks with active
contracts and a complete $\pm 9$-month event window.  CDS spreads directly
price default probability and are traded by sophisticated institutional
counterparties whose mandate is precisely to monitor bank credit quality.
The pre-adoption CDS coefficients cluster tightly around zero through all nine
lead months --- even more precisely than the equity estimates --- confirming
that credit risk monitors did not front-run the CECL disclosure.  At adoption,
spreads jump and widen persistently.  The most precisely estimated post-adoption
coefficient is at month~4: $+190.8$~basis points per unit of CECL-to-equity
($p = 0.047$), implying that a ten-percentage-point higher CECL adjustment is
associated with approximately 19~basis points wider CDS spreads four months
after adoption.  The directional pattern is sustained through month~6.

\paragraph{Joint interpretation.}
The equity and CDS results jointly establish two properties required for our
identification.  First, the CECL adjustment is genuinely informative: even
market participants with strong incentives to already price bank credit risk
revised their assessments sharply and persistently at the moment of disclosure,
ruling out the possibility that CECL was merely a relabeling of already-known
information.  Second, the information was unanticipated: the flat pre-adoption
leads across nine months in both markets confirm no systematic anticipation of
the adjustment's magnitude.  Together, these properties support the
interpretation that when community bank uninsured depositors observe a large
CECL adjustment in their bank's 2023Q1 Call Report, they are receiving a
genuine, previously unavailable signal about the credit-risk composition of
the bank's loan portfolio --- and that subsequent changes in their deposit
behavior reflect a causal response to that signal.


%% ============================================================
%%  SUGGESTED NEW BIB ENTRIES (add to depositor-discipline-references.bib)
%%
%%  These methodology citations are standard in DiD papers at top journals.
%%
%%  @article{callaway2021difference,
%%    title={Difference-in-differences with multiple time periods},
%%    author={Callaway, Brantly and Sant'Anna, Pedro H.C.},
%%    journal={Journal of Econometrics},
%%    volume={225},
%%    number={2},
%%    pages={200--230},
%%    year={2021},
%%    publisher={Elsevier}
%%  }
%%
%%  @article{sun2021estimating,
%%    title={Estimating dynamic treatment effects in event studies with
%%           heterogeneous treatment effects},
%%    author={Sun, Liyang and Abraham, Sarah},
%%    journal={Journal of Econometrics},
%%    volume={225},
%%    number={2},
%%    pages={175--199},
%%    year={2021},
%%    publisher={Elsevier}
%%  }
%%
%%  @article{dechaisemartin2020two,
%%    title={Two-way fixed effects estimators with heterogeneous treatment effects},
%%    author={de Chaisemartin, Cl{\'e}ment and D'Haultf{\oe}uille, Xavier},
%%    journal={American Economic Review},
%%    volume={110},
%%    number={9},
%%    pages={2964--2996},
%%    year={2020},
%%    publisher={American Economic Association}
%%  }
%%
%%  @article{roth2023whats,
%%    title={What's trending in difference-in-differences? {A} synthesis of the
%%           recent econometrics literature},
%%    author={Roth, Jonathan and Sant'Anna, Pedro H.C. and Bilinski, Alyssa
%%            and Poe, John},
%%    journal={Journal of Econometrics},
%%    volume={235},
%%    number={2},
%%    pages={2218--2244},
%%    year={2023},
%%    publisher={Elsevier}
%%  }
%%
%% ============================================================
